%! Author = Sujal Singh
%! Date = 9/25/23

% Preamble
\documentclass[11pt,twocolumn,oneside]{book}

\title{
    \includegraphics{./images/ipu}
    \\
    \textbf{\\Proposal for the Establishment of CSSAC: \\ Computer Science \& Systems Administration Club}
}
\author{Sujal Singh \\ University School of Automation \& Robotics}
\date{Dated: September 25, 2023}

% Packages
\usepackage{amsmath}
\usepackage{lipsum}
\usepackage{lmodern}
\usepackage[T1]{fontenc}
\usepackage[utf8]{inputenc}
\usepackage[a4paper,left=2cm,right=2cm,top=3cm,bottom=3cm]{geometry}
\usepackage[
    pdftitle={CSSAC Proposal},
    pdfsubject={CSSAC Proposal},
    pdfauthor={Sujal Singh},
    pdfdisplaydoctitle,
    hidelinks,
]{hyperref}
\usepackage{color}
\usepackage{graphicx}
\graphicspath{{images/}}

% Document
\begin{document}
    \maketitle
    \tableofcontents


    \chapter{Introduction: The What and Why}\label{ch:introduction}


    \section{Abstract}\label{sec:abstract}
    This proposal requests the establishment of a Computer Science \& Systems Administration Club (CSSAC), an
    exceptionally unique and innovative endeavor, as the existence of such clubs in Indian universities is rare, if not
    non-existent.
    Such a club offers unparalleled learning experiences to members, has practical usefulness to virtually all other
    clubs and has the potential to offer facilities typically only found in the Indian Institutes of Technology and
    various leading universities abroad such as MIT, Stanford, University of California and The University of Waterloo.
    \\
    All details necessary for the creation of CSSAC are outlined in this proposal, including requirements from the
    university, and it also demonstrates that the author is capable of leading the technical aspects of this initiative.


    \section{Vision}\label{sec:vision}
    \emph{``What I cannot create, I do not understand.''}{\par\hfill\small{---Richard Feynman}}
    \\\\
    At the very heart of this club is the goal to learn by doing, to obtain in-depth knowledge, to cultivate a community
    of technically competent members who possess a strong theoretical foundation in computer science but also have the
    ability to practically apply that knowledge to real-world problems.\ After all, we're soon to be engineers.
    CSSAC places no emphasis on tasks aimed solely at securing better job prospects, but only on the goal of
    utilizing our time at the university to become better engineers, and yet it'll allow members to distinguish
    themselves from others in a pool of equally talented candidates.\ CSSAC also strives to promote
    \href{https://wikipedia.org/wiki/Hacker_culture}{\color{blue}\underline{hacker culture}} (not to be confused with
    cracker) throughout the university.
    In practical terms, this translates into having two subteams (not necessarily mutually exclusive), as implied by the
    name of the club.\ These subteams will undertake two different sets of problems described in their respective
    sections, a computer science team dealing with mostly abstract problems and the systems administration team applying
    the learnings from the aforementioned team to practical purposes with fruitful outcomes.

    \section{\large Systems Administration Team}\label{sec:systems-administration-team}
    This systems administration team focuses on addressing practical problems.\ It's primary objectives are as follows:
    \begin{enumerate}
        \item \textbf{On-premise Cloud:} The systems administration team will design, establish and maintain an
        on-premise private cloud offering available as a service to members, other clubs, faculty and individual
        students.\ This also lays a foundation for significantly more ambitious projects in the future such as providing
        high-performance computing (HPC) to research projects in the university.
        \item \textbf{Services:} CSSAC will reserve a portion of the resources to deploy services for the
        entire\footnote{Limited to those who register for an account and available capacity.} university.\ This will
        include services like providing cloud storage, git hosting, mail services and mailing lists, static web hosting,
        etc.\ Initially, we might reserve all resources to deploy services that are available to everyone in the
        university.
    \end{enumerate}
    The exact implementation details are not discussed in this proposal since this will mostly be dictated by the
    hardware provided to the club.

    \subsection{Symbiotic Relationship with Other Clubs}\label{subsec:symbiotic-relationship-with-other-clubs}
    This club has the potential to offer significant utility to other clubs and offer them the opportunity to develop
    various innovative projects in collaboration such that these projects are mutually beneficial to all parties
    involved.\ Several such scenarios are described in the following sections:

    \subsubsection{IoSC, GDSC, SDC, ACM and IEEE}
    Many clubs in our university have various teams, for example, IoSC has teams such as the software development team
    (i3), IoT and cloud team (i7) and AI development team (i9), all of which stand to gain immense benefits from the
    compute services offered by CSSAC.\ We could also offer GPU compute services in the future to be able to train
    machine learning models and other applications which heavily benefit from GPU compute.
    \\
    Students who are particularly interested in cloud computing would be thrilled to learn more about how this club
    plans to set up a mini data center on campus and if they choose to join the systems administration team of our club
    then they'll gain invaluable hands-on experience.
    \\
    CSSAC also stands to benefit from these clubs; for example, instead of creating our own web development team we
    could offload such tasks to the web development teams of these clubs\@.

    \subsubsection{Karuyantra Club}
    The Karuyantra club could offer us 3D printing services to be able to produce various makeshift solutions in an
    example scenario such as when a network card's bracket will not fit the chassis of the server.\ They can also
    work on a project to develop a robot which can remotely perform tasks in a data center that would normally require a
    human's physical presence such as replacing faulty hard drives or periodic reapplication of thermal paste on CPUs.

\end{document}